% Generated from locale/tr/content/intuitionistic-logic/tableaux/rules.tex; do not hand-edit.


\olfileid[tr]{int}{tab}{rul}

\olsection{Sezgisel Mantık için Kurallar}

$\land$ ve $\lor$ bağlaçlarının kuralları, yalnızca önekler eklenmiş olmak
üzere, olağan önermesel işaretli !!{tableau}s kurallarıyla aynıdır. Her
durumda, işaretli !!{formula}~$\sFmla{S}{!A}[\sigma]$'ya uygulanan kural,
yine $\sigma$ ile öneklenmiş yeni !!{formula}s üretir. Bu sezgisel olarak
açıktır: örneğin $!A \land !B$, $\sigma$'nın adlandırdığı dünyada doğruysa
$!A$ ile $!B$ de $\sigma$'da (başka bir dünyada değil) doğrudur. $\land$ ve
$\lor$ kurallarını \olref{tab:prop-rules}'te bir araya getiriyoruz.

\begin{table}
  \[\def\arraystretch{3}\begin{array}{|c|c|}
    \hline
    \AxiomC{\sFmla{\True}{!A \land !B}[\sigma]}
    \RightLabel{\TRule{\True}{\land}}
    \UnaryInfC{\sFmla{\True}{!A}[\sigma]}
    \noLine
    \UnaryInfC{\sFmla{\True}{!B}[\sigma]}
    \DisplayProof
    &
    \AxiomC{\sFmla{\False}{!A \land !B}[\sigma]}
    \RightLabel{\TRule{\False}{\land}}
    \UnaryInfC{$\sFmla{\False}{!A}[\sigma] \quad \mid \quad
      \sFmla{\False}{!B}[\sigma]$}
    \DisplayProof
    \\[2ex]
    \hline
    \AxiomC{\sFmla{\True}{!A \lor !B}[\sigma]}
    \RightLabel{\TRule{\True}{\lor}}
    \UnaryInfC{$\sFmla{\True}{!A}[\sigma] \quad \mid \quad
      \sFmla{\True}{!B}[\sigma]$}
    \DisplayProof
    &
    \AxiomC{\sFmla{\False}{!A \lor !B}[\sigma]}
    \RightLabel{\TRule{\False}{\lor}}
    \UnaryInfC{\sFmla{\False}{!A}[\sigma]}
    \noLine
    \UnaryInfC{\sFmla{\False}{!B}[\sigma]}
    \DisplayProof
    \\[2ex]
    \hline
  \end{array}\]
  \caption{$\land$ ve $\lor$ için önekli !!{tableau} kuralları}
  \ollabel{tab:prop-rules}
\end{table}

Kapanma koşulu olağan !!{tableau}s için olana benzer; ancak yalnızca
!!{formula}s{}in değil, öneklerin de eşleşmesini isteriz. Aslında biraz daha
esnek olabiliriz: Sezgisel modellerde !!{formula}s bir kez doğru olduktan
sonra doğru kalır; dolayısıyla $!A$'nın $\sigma$'da doğru, fakat erişilebilir
bir $\sigma.{*}$ önekinde yanlış olması olanaksızdır. Bu nedenle bir dal,
bazı $\sigma$ öneki ve !!{formula}~$!A$ için hem
\[
\sFmla{\True}{!A}[\sigma] \quad\text{hem de}\quad
\sFmla{\False}{!A}[\sigma.{*}]
\]
içeriyorsa kapalıdır. İşaretler ters çevrilmişse, yani dal
\[
\sFmla{\False}{!A}[\sigma] \quad\text{ve}\quad
\sFmla{\True}{!A}[\sigma.{*}]
\]
içeriyorsa yalnızca $*$ boş dizi olduğunda kapalıdır.

Ayrıca bir dal $\sFmla{\True}{\bot}[\sigma]$ içeriyorsa kapalıdır.

Varsayımları yerleştirme kuralı da olağan !!{tableau}s için olanla aynıdır;
tek fark, varsayımlar için her zaman $1$ önekini kullanmamızdır. (Bütün
varsayımlarda aynı öneki kullandığımız sürece hangisini seçtiğimiz önemli
değildir.) Örneğin
\[
!B_1, \dots, !B_n \Proves !A
\]
ancak ve ancak şu varsayımlar için kapalı bir tablo varsa geçerlidir:
\[
\sFmla{\True}{!B_1}[1], \dots, \sFmla{\True}{!B_n}[1],
\sFmla{\False}{!A}[1].
\]

Koşullu $\lif$ için kurallar klasik ve modal durumlardakinden farklıdır.
$\TRule{\True}{\lif}$ kuralı,
$\sFmla{\True}{!A \lif !B}[\sigma]$ içeren bir dalı farklı iki dal üzerinde
$\sFmla{\False}{!A}[\sigma.{*}]$ ve
$\sFmla{\True}{!B}[\sigma.{*}]$ ile genişletir. Bu kural yalnızca uygulandığı
dalda \emph{zaten} bulunan bir $\sigma.{*}$ öneki için uygulanabilir. Böyle
bir öneke dalda ``kullanılmış'' diyelim. ($\sigma.{*}$, $\sigma$'nın kendisini
de içerdiğinden, kural her zaman $\sFmla{\False}{!A}[\sigma]$ ve
$\sFmla{\True}{!B}[\sigma]$ önekli işaretli formüllerini ayrı dallara
ekleyerek uygulanabilir.)

$\TRule{\False}{\lif}$ kuralı,
$\sFmla{\False}{!A \lif !B}[\sigma]$ içeren bir dalı aynı dal üzerinde hem
$\sFmla{\True}{!A}[\sigma.n]$ hem de $\sFmla{\False}{!B}[\sigma.n]$ ile
genişletir; burada $\sigma.n$ dalda yeni bir önektir.

$\lnot$ kuralları da benzer biçimde ($\lnot !A$'nın
$!A \lif \lfalse$ tanımı kullanılarak) tanımlanır.

Kurallar \olref{tab:rules-lif-lnot}'ta verilmiştir.

\begin{table}
  \[\def\arraystretch{3}\begin{array}{|c|c|}
    \hline
    \AxiomC{\sFmla{\True}{\lnot !A}[\sigma]}
    \RightLabel{\TRule{\True}{\lnot}}
    \UnaryInfC{$\sFmla{\False}{!A}[\sigma.{*}]$}
    \DisplayProof
    &
    \AxiomC{\sFmla{\False}{\lnot !A}[\sigma]}
    \RightLabel{\TRule{\False}{\lnot}}
    \UnaryInfC{\sFmla{\True}{!A}[\sigma.n]}
    \DisplayProof
    \\[1ex]
    \text{$\sigma.{*}$ kullanılmıştır} & \text{$\sigma.n$ yenidir}\\
    \hline
    \AxiomC{\sFmla{\True}{!A \lif !B}[\sigma]}
    \RightLabel{\TRule{\True}{\lif}}
    \UnaryInfC{$\sFmla{\False}{!A}[\sigma.{*}] \quad \mid \quad
      \sFmla{\True}{!B}[\sigma.{*}]$}
    \DisplayProof
    &
    \AxiomC{\sFmla{\False}{!A \lif !B}[\sigma]}
    \RightLabel{\TRule{\False}{\lif}}
    \UnaryInfC{\sFmla{\True}{!A}[\sigma.n]}
    \noLine
    \UnaryInfC{\sFmla{\False}{!B}[\sigma.n]}
    \DisplayProof
    \\[1ex]
    \text{$\sigma.{*}$ kullanılmıştır} & \text{$\sigma.n$ yenidir}\\
    \hline
  \end{array}\]
  \caption{$\lnot$ ve $\lif$ için önekli !!{tableau} kuralları}
  \ollabel{tab:rules-lif-lnot}
\end{table}

